\section{Introduction}

Generative modeling aims to transform a simple prior distribution into a complex target data distribution. Among existing approaches, flow-based models provide an understandable method by explicitly defining an invertible transformation between random variables and computing the corresponding Jacobian determinants. In contrast, the Flow Matching framework~\cite{flowmatching} achieves this by modeling an instantaneous velocity field that transports one distribution into another based on the continuity equation. Both approaches are built on concepts of ``flows'' and construct an explicit mapping from the prior to the target distribution. However, they differ in what the transform is and how it is learned.

Recently, research has paid more attention to few-step, and particularly, one-step generative models which offer significant sampling advantages. The MeanFlow model~\cite{meanflow} extends Flow Matching by introducing a ground-truth average velocity field derived from the instantaneous velocity field in Flow Matching. The average velocity is defined as the ratio of displacement to a time interval, with the displacement computed as the integral of the instantaneous velocity. This formulation allows the neural network to directly learn from the well-defined average velocity without any extra consistency heuristic~\cite{meanflow}, enabling reliable and effective one-step generation. Moreover, classifier-free guidance (CFG) can be incorporated into the target vector field without additional cost at sampling time, preserving model's efficiency.

In this work, we focus on Flow Matching and MeanFlow. Our primary goal is to understand their fundamental principles through theoretical derivation and hands-on experiments. Building on this foundation, we adapt the MeanFlow framework for one-step image denoising task. Specifically, we alters its initial distribution to noisy images instead of gaussian noise, retaining its one-step generation ability and serving the denoising purpose. Experiments on MNIST demonstrate that MeanFlow denoising method can substantially reduce noise with just a single function evaluation, highlighting its potential as an efficient one-step denoising method.
